\documentclass{article}

\textwidth=6.5in
\textheight=8.5in
\voffset=-54pt
\oddsidemargin=0in
\evensidemargin=0in
\setlength{\parskip}{8pt}
\setlength{\parindent}{0pt}

\usepackage{workbook}

\usepackage[sols]{handout}

\begin{document}

\begin{center}
  \large\textsc{Problem Set 34 - Introduction to Differential Equations}
  
      \pspace{12pt}
  \begin{problemonly}\large Name: \rule{5cm}{0.15mm}\\ \end{problemonly}
\end{center}

\begin{grading}
  
    \begin{enumerate}
    \item 2 points
    \item 12 points 
    \item 5 points
    \item 3 points - 2/1
    \item 5 points - 2/3
    \item 2 points 
    \end{enumerate}
  Total: 29 points
\end{grading}

\begin{enumerate}

\item 



 %(variation on 31.2.1) % 31.2.1

  \begin{problem}
    Which two of the following are solutions to the differential equation
    $y'(t)=-5y$?  Be sure to pick two solutions.
  \end{problem}

  
    \begin{enumerate}[label=\protect\maybecircle{\Alph*.}]
    \begin{multicols}{4}
    \plainitem $y=e^{5t}$
    \plainitem $y= - \frac{5}{2} t^2$
    \circleitem $y=e^{-5t}$
    \plainitem $y=\sin(5t)$
    \plainitem $y=e^{5t} + 7$
    \plainitem  $y=7e^{5t} $
    \plainitem $y=e^{-5t} + 3$
    \circleitem $y=3e^{-5t} $
    \end{multicols}
    \end{enumerate}    
  

  \begin{solution}
    Let's check each one:
    \begin{enumerate}[label=\protect\maybecircle{\Alph*.}]      
    \item If $y(t) = e^{5t}$, then differentiating gives us $y'(t) = 5
      e^{5t}$.  On the other hand, $-5 y(t) = -5 e^{5t}$, so $y'(t)$ and
      $-5 y(t)$ are not the same.
    \item If $y(t) = - \frac{5}{2} t^2$, then $y'(t) = -5t$ while $- 5
      y(t) = -\frac{25}{2} t^2$; so, $y'(t)$ and $-5 y(t)$ are not the
      same.
    \item If $y(t) = e^{-5t}$, then $y'(t) = - 5 e^{-5t}$ while $-5 y(t) =
      - 5 e^{-5t}$, so $y'(t)$ and $-5 y(t)$ are the same in this case.
    \item If $y(t) = \sin(5t)$, then $y'(t) = 5 \cos(5t)$ while $-5 y(t) =
      -5 \sin(5t)$; these are not the same.
    \item If $y(t) = e^{5t} + 7$, then $y'(t) = 5 e^{5t}$ while $-5 y(t) = -5
      (e^{5t} + 7) = - 5 e^{5t} - 35$; these are not the same.
    \item If $y(t) = 7 e^{5t}$, then $y'(t) = 35 e^{5t}$ while $-5 y(t) =
      -35 e^{5t}$; these are not the same.
    \item If $y(t) = e^{-5t} + 3$, then $y'(t) = - 5 e^{-5t}$ and $-5 y(t)
      = -5 \left( e^{-5t} + 3 \right) = -5 e^{-5t} - 15$; these are not
      the same.
    \item If $y(t) = 3 e^{-5t}$, then $y'(t) = -15 e^{-5t}$ and $-5 y(t) =
      -15 e^{-5t}$; these are equal. 
    \end{enumerate}
  \end{solution}

\pspace{\vfill\newpage}

\item 

\begin{grading}
    12 points. 3 for each differential equation. For each solution they pick, they must
show that the left-hand side equals the right-hand side. If they don’t, they should receive little or no
credit.
\end{grading}

 % 31.2.5

  \begin{problem}
    For help with this problem read the two examples on
    pp. 991-992 in Gottlieb.

    Determine which of the following functions are solutions to each of
    the differential equations below. A given differential equation may
    have more than one solution.
    
    \emph{Recall that $\displaystyle \frac{d^2 y}{dt^2}$ means $\displaystyle \frac{d}{dt}\bigg(\frac{dy}{dt}\bigg)$, the second derivative of $y$ with respect to $t$.}
  \end{problem}

      Differential equations:
    
      \begin{enumerate}
      \begin{multicols}{4}
      \item $\displaystyle \frac{dy}{dt} = t$
      \item $\displaystyle \frac{dy}{dt}=y$
      \item $\displaystyle \frac{dy}{dt}=e^t$
      \item$\displaystyle \frac{d^2 y}{dt^2}=4y$
      \end{multicols}
      \end{enumerate}
  
    Solution choices:    
      \begin{enumerate}[label=\roman*., ref=\roman*]
      \begin{multicols}{4}
      \item $y=\frac{t^2}{2}$
\item $y=\frac{t^2}{2}+5$
\item $y=e^{-2t}$
\item $y=e^t+5$
\item $y=2e^t$
\item $y=e^{2t}$
\item $y=5e^{2t}$
\item $y=e^{2t}+5$
\end{multicols}
\end{enumerate}

      \pspace{\vfill}


  \begin{solution}
    \begin{enumerate}
    \item The solutions of $\displaystyle \frac{dy}{dt} = t$ are
      {i} and {ii}.
      This differential equation is saying that the derivative of $y(t)$
      is the function $t$, so $y(t)$ is an antiderivative of $t$; that is,
      $y(t) = \displaystyle \int t\,dt = \frac{t^2}{2} + C$.

    \item Among the choices, the only one that is a solution of
      $\displaystyle \frac{dy}{dt}=y$ is $y = 2 e^t$, choice
      {v}.  After all, we know that the general
      solution of $\frac{dy}{dt} = y$ is $y = C e^t$. 

    \item Among the choices, the only one that is a solution of
      $\displaystyle \frac{dy}{dt} = e^t$ is $y = e^t + 5$, choice
      {iv}.  This differential equation is saying
      that $y$ is an antiderivative of $e^t$, so the general solution is
      $y(t) = \displaystyle \int e^t\,dt = e^t + C$.

    \item Choices {iii},
      {vi}, and {vii} are
       solutions of $\displaystyle \frac{d^2 y}{dt^2}=4y$.  (We don't know
       the general solution of this differential equation, but we can
       check each choice.)

    \end{enumerate}
  \end{solution}

\pspace{\newpage}

\item 

\begin{grading}
    5 points. They need to check all 5 choices.
\end{grading}

  \begin{problem}
    Which one of the following is a solution to the differential equation
    $\displaystyle \frac{dy}{dt}=y+1$?
  \end{problem}

  
    \begin{enumerate}[label=\protect\maybecircle{\Alph*.}]
    \plainitem $y=Ce^{t}$
    \plainitem $y=Ce^t-t$
    \plainitem $y=C(t^2+t)$
    \circleitem $y=Ce^t-1$
    \plainitem $y=Ce^{-t}+1$
    \end{enumerate}
  

  \begin{solution}
    We just check each one:
    \begin{enumerate}[label=\Alph*.]
    \item If $y = C e^t$, then $\frac{dy}{dt} = C e^t$ while $y + 1 = C
      e^t + 1$, so $\frac{dy}{dt} \neq y + 1$.
    \item If $y = C e^t - t$, then $\frac{dy}{dt} = C e^t - 1$ while $y +
      1 = C e^t - t + 1$, so $\frac{dy}{dt} \neq y + 1$.
    \item If $y = C (t^2 + t)$, then $\frac{dy}{dt} = C (2t + 1)$ while $y
      + 1 = C (t^2 + t) + 1$, so $\frac{dy}{dt} \neq y + 1$.
    \item If $y = C e^t - 1$, then $\frac{dy}{dt} = C e^t$ while $y + 1 =
      C e^t$, so $\frac{dy}{dt} = y + 1$.
    \item If $y = C e^{-t} + 1$, then $\frac{dy}{dt} = - C e^{-t}$ while
      $y + 1 = C e^{-t} + 2$, so $\frac{dy}{dt} \neq y + 1$. 
    \end{enumerate}
  \end{solution}

\pspace{\vfill}

\item

% Meghan / Meredith?

  The differential equations $\frac{dy}{dt} = t^3$ and $\frac{dy}{dt} =
  y^3$ look similar but are in fact very different.

  \begin{enumerate}
  \item
\begin{problem}[\vfill]
      It should be relatively easy to solve one of these two. Write the
      general solution for the easier equation and explain in one complete
      sentence what your method was.
    \end{problem}

    \begin{solution}
      It's easy to solve $\frac{dy}{dt} = t^3$: this equation says that
      $y$ is an antiderivative of $t^3$, so the general solution is $y(t) =
      \displaystyle \int t^3\,dt = \boxed{\frac{t^4}{4} + C}$.
      So, our method was to find all antiderivatives of $t^3$. 
    \end{solution}

  \item
    \begin{problem}[\vfill\newpage]
      The more challenging equation can't be solved using the method you
      used in {{{{\#4}}(a)}}. Why is that? What is the key
      difference?
    \end{problem}

    \begin{solution}
      The differential equation $\frac{dy}{dt} = y^3$ says that $y$ is an
      antiderivative of $y^3$; the problem is that $y$ is a function that
      we don't know yet, so we don't have a concrete function to
      anti-differentiate.
    \end{solution}

  \end{enumerate}

\item 

%(Almost identical to 31.2.12) 

  %($\star$) 
  When a population has unlimited resources and is free from
  disease and strife, the rate at which the population grows in often
  modeled as being proportional to the population size.  That is, if
  $P(t)$ is the population size at time $t$, then $\dfrac{dP}{dt} = kP$
  where $k$ is a constant.  (The constant $k$ is called the
  \uline{relative growth rate} of the population.)  Assume that both the
  bug and the minnow populations described below behave according to this
  model.

  In both scenarios described below you are given enough information to
  find the relative growth rate $k$. In one case, you can find $k$ using
  just the differential equation; in the second scenario, you must
  actually use the solutions of the differential equation to find $k$.

  \begin{enumerate}

  \item
    \begin{problem}[\vfill]
      Let $M=M(t)$ be the minnow population at time $t$, $t$ in weeks. At
      $t=0$ there are 500 minnows.  When there are 700 minnows, the
      population is growing at a rate of 50 minnows per week.  Write a
      differential equation reflecting the situation.  (Note: We've
      already said above that we want to model this as $\dfrac{dM}{dt} =
      kM$, so you just need to find $k$.)
    \end{problem}

    \begin{solution}
      We are told that, when $M = 700$, $\frac{dM}{dt} = 50$; we can plug
      this information directly into the differential equation
      $\frac{dM}{dt} = kM$ to solve for $k$: $50 = k (700)$, so $k =
      \frac{1}{14}$.  So, the differential equation is
      $\boxed{\frac{dM}{dt} = \frac{1}{14}M}$.
    \end{solution}

  \item
\begin{problem}[\vfill]
      Let $B=B(t)$ be the bug population at time $t$, $t$ in weeks. At
      $t=0$ there are 6000 bugs. When $t=5$ there are 9000 bugs. Write a
      differential equation reflecting the situation.
    \end{problem}

    \begin{solution}
      We're modeling this as $\frac{dB}{dt} = kB$, but we are not told
      $\frac{dB}{dt}$ at any time, so we need to solve the differential
      equation in order to determine $k$.

      The general solution of $\frac{dB}{dt} = kB$ is $B(t) = C e^{kt}$.
      We're told that $B(0) = 6000$ and $B(5) = 9000$; let's see what
      these tell us:
      \begin{align*}
        B(0) &= 6000 \Longrightarrow C e^{k \cdot 0} = 6000
        \Longrightarrow C = 6000 \\
        B(5) &= 9000 \Longrightarrow C e^{5k} = 9000
      \end{align*}
      Since we already figured out that $C = 6000$, the second equation
      tells us that
      \begin{align*}
        6000 e^{5k} &= 9000 \\
        e^{5k} &= \frac{9000}{6000} = \frac{3}{2} \\
        5k &= \ln \left( \frac{3}{2} \right) \\
        k &= \frac{1}{5} \ln \left( \frac{3}{2} \right)
      \end{align*}
      So, our differential equation is $\boxed{\frac{dB}{dt} = \frac{1}{5}
        \ln \left( \frac{3}{2} \right) B}$. 
    \end{solution}

  \end{enumerate}



\pspace{\vfill\newpage}

\item
\begin{enumerate}
  \item
    \begin{problem}[\vfill]
      Yoshi currently has $Y$ dollars in his bank account.  If the account has a
      nominal annual interest rate of 7\%, compounded continuously, how
      much money will Yoshi have in $t$ years?  (Assume Yoshi makes no
      deposits or withdrawals.)
    \end{problem}

   \emph{If you'd like to refresh your understanding
         of compound interest, see
          \#4 in section 4.3 of the workbook and
          Problem Set 19, \#2. You may use the continuous compounding formula, $A=Pe^{rt}$.}

    \begin{solution}
      This is continuous compounding, so the amount of money Yoshi will have in $t$ years is $Y e^{0.07 t}$.
    \end{solution}

  \item
\begin{problem}[\vfill]
      Let $M(t)$ be the amount of money that Yoshi has in $t$ years.
      Which of the following differential equations is $M(t)$ a solution
      of?
      
        \begin{enumerate}[label=\protect\maybecircle{\Alph*.}]
        \plainitem $\frac{dM}{dt} = 0.07$
        \plainitem $\frac{dM}{dt} = 1.07$
        \plainitem $\frac{dM}{dt} = 0.07 t$
        \circleitem $\frac{dM}{dt} = 0.07 M$
        \end{enumerate}
      
    \end{problem}

    \begin{solution}
      The derivative of $M(t) = Y e^{0.07 t}$ is $\frac{dM}{dt} = 0.07 Y
      e^{0.07 t}$, which is $0.07 M(t)$. 
    \end{solution}

  % \item
  %   \begin{problem}[\vfill]
  %     Interpret the differential equation you chose in
  %     words.  That is, if you had to describe to a friend what the
  %     differential equation says about how fast Yoshi's money is growing, what would you say?
  %   \end{problem}

  %   \begin{solution}
  %     The differential equation $\frac{dM}{dt} = 0.07 M$ can be read as
  %     saying, ``The (instantaneous) rate at which the account earns
  %     interest is always equal to 7\% of the amount of money.''  
  %   \end{solution}

  % \item
  %   \begin{problem}[\vfill]
  %     How long does it take Yoshi's money to double?
  %   \end{problem}

  %   \begin{solution}
  %     If Yoshi starts with $\$Y$, we'd like to know when his amount of
  %     money reaches $\$2Y$.  That is, we'd like to solve
  %     \begin{align*}
  %       Y e^{0.07 t} &= 2Y \\
  %       e^{0.07 t} &= 2 \\
  %       0.07 t &= \ln 2 \\
  %       t &= \frac{\ln 2}{0.07}
  %     \end{align*}
  %     So, it takes $\boxed{\frac{\ln 2}{0.07} \approx 9.9}$ years for
  %     Yoshi's money to double. 
  %   \end{solution}

  \end{enumerate}

\item 

 \textbf{Reflect Back.}
You know
  enough to find the general solution of one of the following differential
  equations.  Which one?  Find its general solution.

  
    \begin{enumerate}
    \item $\dfrac{dy}{dt} = ty + 1$
    \item $\dfrac{dy}{dt} = y^2 + 1$
    \item
$\dfrac{dy}{dt} = t^2 + 1$
    \end{enumerate}    

    \pspace{\vfill}
  

  \begin{solution}
    The one that is simple is {{(c)}}: it
    says that $y$ is an antiderivative of $t^2 + 1$, so the general
    solution is $y = \displaystyle \int (t^2 + 1)\,dt = \boxed{\frac{1}{3}
      t^3 + t + C}$.
  \end{solution}  

\end{enumerate}
\spaceonly{\psreflection{
Optional reading for next class is \S31.1 in Gottlieb.
}}

\end{document}